\begin{tikzpicture}[
  x=0.94cm,y=0.90cm,
  wquvalue/.style={font=\scriptsize,fill=white,inner sep=0.65pt,
    text=black!82},
  wqmvalue/.style={font=\tiny,fill=white,inner sep=0.45pt,
    text=black!82}
]
\path[use as bounding box] (-0.20,-0.28) rectangle (16.15,5.30);
\draw[wqstring] (0,2.00)--(16,2.00);
\draw[wqstring] (0,0.40)--(16,0.40);
\node[wqlevel,anchor=east] at (-0.12,2.00) {$1$};
\node[wqlevel,anchor=east] at (-0.12,0.40) {$2$};

\node[wqlargefrozen]  (q10) at (1.20,2.00) {$[0,1]$};
\node[wqlargemutable] (q11) at (3.60,2.00) {$[1,2]$};
\node[wqlargemutable] (q12) at (7.30,2.00) {$[2,5]$};
\node[wqlargemutable] (q13) at (11.00,2.00) {$[5,6]$};
\node[wqlargefrozen]  (q14) at (14.10,2.00) {$[6,\infty]$};
\node[wqlargefrozen]  (q20) at (3.20,0.40) {$[0,3]$};
\node[wqlargemutable] (q21) at (7.25,0.40) {$[3,4]$};
\node[wqlargemutable] (q22) at (11.10,0.40) {$[4,7]$};
\node[wqlargefrozen]  (q23) at (15.05,0.40) {$[7,\infty]$};

% The two additional frozen vertices indexed by T.
% Their horizontal positions balance the string vertices incident to them.
\node[wqweavefrozen] (ext-sigma) at (5.45,4.05) {};
\node[wqweavefrozen] (ext-tau) at (9.05,4.05) {};
\node[wqvlabel,above=0.7mm of ext-sigma] {$\sigma$};
\node[wqvlabel,above=0.7mm of ext-tau] {$\tau$};

% Mixed blocks: each nonzero entry a_rho(t)=1 contributes the path
% left_beta(t) -> rho -> right_beta(t).
\wqextensiontriangle{q11}{ext-sigma}{q12}
\wqextensiontriangle{q11}{ext-tau}{q12}
\wqextensiontriangle{q21}{ext-tau}{q22}
\wqextensiontriangle{q13}{ext-tau}{q14}

% Draw the string-quiver block above the mixed block, as in the local
% extended-mutation figure.  The dotted arrows are the half-weight edges.
\draw[wqstringarrow] (q11)--(q10);
\draw[wqstringarrow] (q12)--(q11);
\draw[wqstringarrow] (q13)--(q12);
\draw[wqstringarrow] (q14)--(q13);
\draw[wqstringarrow] (q21)--(q20);
\draw[wqstringarrow] (q22)--(q21);
\draw[wqstringarrow] (q23)--(q22);
\draw[wqstringhalf] (q10)--(q20);
\draw[wqstringarrow] (q20)--(q12);
\draw[wqstringarrow] (q12)--(q22);
\draw[wqstringarrow] (q22)--(q14);
\draw[wqstringhalf] (q14)--(q23);

% Cartan monomials, placed immediately below their string vertices.
\node[wqmvalue,anchor=north] at ([yshift=-0.65mm]q10.south) {$1$};
\node[wqmvalue,anchor=north] at ([yshift=-0.65mm]q11.south) {$1$};
\node[wqmvalue,anchor=north] at ([yshift=-0.65mm]q12.south) {$A_\sigma A_\tau$};
\node[wqmvalue,anchor=north] at ([yshift=-0.65mm]q13.south) {$A_\sigma^{-1}$};
\node[wqmvalue,anchor=north] at ([yshift=-0.65mm]q14.south) {$A_\sigma A_\tau^2$};
\node[wqmvalue,anchor=north] at ([yshift=-0.65mm]q20.south) {$1$};
\node[wqmvalue,anchor=north] at ([yshift=-0.65mm]q21.south) {$A_\sigma A_\tau$};
\node[wqmvalue,anchor=north] at ([yshift=-0.65mm]q22.south) {$A_\tau$};
\node[wqmvalue,anchor=north] at ([yshift=-0.65mm]q23.south) {$A_\sigma A_\tau$};

% u-monomials, placed at the corresponding word positions.
\node[wquvalue] at (2.40,2.00) {$1$};
\node[wquvalue] at (5.45,2.00) {$A_\sigma A_\tau$};
\node[wquvalue] at (5.23,0.40) {$1$};
\node[wquvalue] at (9.18,0.40) {$A_\tau$};
\node[wquvalue] at (9.15,2.00) {$1$};
\node[wquvalue] at (12.55,2.00) {$A_\tau$};
\node[wquvalue] at (13.08,0.40) {$1$};
\end{tikzpicture}
